
\documentclass[11pt,a4paper]{article}
\usepackage[T1]{fontenc}
\usepackage[utf8]{inputenc}
\usepackage[italian,english]{babel}
\usepackage{lmodern}
\usepackage[margin=2.2cm]{geometry}
\usepackage{amsmath}
\usepackage{enumitem}
\usepackage{booktabs}
\usepackage{titlesec}
\usepackage{microtype}
\usepackage[hidelinks]{hyperref}
\setlength{\parindent}{0pt}
\setlength{\parskip}{0.55em}
\setlength{\emergencystretch}{1.5em}
\clubpenalty=10000
\widowpenalty=10000
\setlist{nosep,leftmargin=*}
\titleformat{\section}{\large\bfseries}{}{0pt}{}
\titlespacing*{\section}{0pt}{1.3ex plus .2ex}{0.6ex}
\newcommand{\topic}[1]{\section*{#1}}
\title{Evolution of the Graphics Processing Unit (GPU)\\[0.3em]
\large Summary}
\author{}
\date{}
\begin{document}
\maketitle
\vspace{-2.7em}
{\small W. J. Dally, S. W. Keckler, and D. B. Kirk.
\emph{IEEE Micro} 41(6), 42--51, 2021.}

\topic{Context}
\begin{otherlanguage}{italian}
Le GPU sono passate da unità finalizzate al rendering 3D per il gaming a GPU
che integrano componenti specifici per l'AI, come i Tensor Core.
L'hardware e il software (tramite CUDA, ad esempio) hanno seguito le esigenze
e il loro cambiamento nel tempo.
\end{otherlanguage}

The central argument is that programmability enables new applications, which
then motivate specialized hardware and software. The examples and performance
figures refer to the period covered by the 2021 article.

\topic{Origin of the first GPU: GeForce 256}
Early 3D accelerators performed fragment operations such as rasterization,
texture mapping, Z-buffering, and shading, while vertex transformation,
clipping, and lighting remained on the CPU. In 1999, \textbf{GeForce 256}
integrated vertex transformation and lighting with fragment processing,
supporting more complex geometry while offloading the CPU. Its pipeline
was still fixed-function. Demand for richer effects subsequently led to
programmable vertex shaders in GeForce 3 and fragment shaders in GeForce FX.

\topic{Programmability and HPC}
Early general-purpose GPU computing (GPGPU) exploited programmable shaders'
floating-point throughput, but required scientific algorithms to be expressed
as rendering tasks. Fragment shaders read textures and interpolated attributes
and wrote to the frame buffer. Stream-processing systems such as Brook-GPU
helped abstract these restrictions and influenced later designs.

\textbf{G80 (Tesla)} and CUDA, launched in 2006, made general-purpose computing
directly accessible. Streaming multiprocessors ran both graphics shaders
and computational threads; shared memory supported communication within
cooperative thread arrays. CUDA provided a C-like, thread- and data-parallel
programming system, removing the need to encode numerical work as graphics.

\textbf{Fermi} strengthened FP64 support and reliability through ECC on main
memory and parity protection for on-chip memories. These features addressed
accuracy and long-running HPC workloads. CUDA Fortran, OpenACC, and libraries
such as cuBLAS, cuSPARSE, and cuFFT reduced porting effort. Together with
education and reusable software, they enabled large scientific systems
such as Titan and Summit.

\topic{Deep learning revolution and Tensor Cores}
Deep learning required algorithms, datasets, and enough computing power to
train networks in practical time. GPUs supplied the computational resource:
AlexNet was trained on two GeForce GTX 580 GPUs in 2012. cuDNN then supplied
optimized neural-network primitives. Unlike many FP64-oriented HPC workloads,
many deep learning operations, especially inference, can use lower precision
to obtain greater throughput, provided accuracy and dynamic range remain adequate.

The article reports a $317\times$ increase in GPU inference performance
from Kepler in 2012 to Ampere in 2020. This trend, called \textbf{Huang's law},
reflects specialized data types and operations beyond transistor scaling
alone. It is not a replacement for Moore's law or equivalent growth in
every FP32 application.

Pascal added FP16 arithmetic and dot-product instructions. Volta introduced
\textbf{Tensor Cores} for mixed-precision matrix multiply-and-accumulate
operations, such as $D=AB+C$, using FP16 multiplication and FP32 accumulation.
Performing many arithmetic operations per specialized instruction amortizes
instruction-fetch, decode, and operand-handling overhead. Turing added integer
Tensor Core operations, including INT8, particularly useful for inference.
This combines a programmable architecture with efficient domain-specific hardware.

\topic{Horizontal scalability}
Larger models require both \emph{scale-up}, improving one GPU, and
\emph{scale-out}, distributing work across GPUs. \textbf{NVLink} and
\textbf{NVSwitch} connect GPUs within nodes, while \textbf{InfiniBand}
connects nodes into clusters. Selene illustrates this hierarchy with 560
eight-GPU nodes, totaling 4,480 GPUs. Efficient communication is essential
to benefit from the additional arithmetic resources.

Software also determines realized performance: TensorRT optimizes inference,
Triton manages serving, and CUTLASS supports matrix kernels. The article
reports training benchmark improvements of up to $2.1\times$ on unchanged
hardware from software enhancements alone.

\topic{Ampere A100 as an architectural case study}
The A100 extends Tensor Cores to formats including BFloat16, TF32, and FP64.
TF32 retains FP32's exponent range with reduced significand precision.
\textbf{Structured sparsity}, requiring two zeros per group of four weights,
can double effective throughput by skipping supported zero-valued work.
This requires the appropriate pattern and model preparation; it is not an
automatic benefit for arbitrary sparse matrices.

The distinction between dense and sparse rates matters: A100's FP16 rates
are 312 TFLOPS dense and 624 TFLOPS structured-sparse, compared with V100's
125 TFLOPS dense. A100 also improves bandwidth and cache capacity, direct
transfers into shared memory, cache residency control, compression, and
task synchronization. Separately, the Ampere graphics family integrates
RT Cores for ray tracing and uses Tensor Cores for DLSS; these graphics
features are not features of A100 itself.

\topic{Memory bottlenecks and the road ahead}
Many HPC applications need approximately 1--4 bytes from main memory per
FP64 operation, while A100 provides about 0.1 byte per peak FP64 operation.
Without sufficient reuse, memory bandwidth therefore limits performance.
The authors identify improved memory technology and packaging, larger
clusters, domain-specific acceleration, and better programming abstractions
as future directions. Their 2021 outlook also emphasizes nonuniform access
costs and placing computation near its data.

\topic{Main contribution and perspective}
The contribution is a historical explanation of how programmability,
specialization, data movement, and software jointly drove GPU adoption.
Its scope is predominantly NVIDIA's evolution, so its examples are not
a comparative survey of all GPU vendors.

\end{document}
